% 使用 ctexart 文档类，支持中文
\documentclass[12pt, a4paper]{ctexart}

% ========== 导言区：引入宏包并进行设置 ==========
\usepackage{geometry} % 用于设置页边距
\geometry{left=2.5cm, right=2.5cm, top=2.5cm, bottom=2.5cm}

\usepackage{graphicx} % 用于插入图片
\usepackage{float}    % 用于控制图片位置 [H]
\usepackage{caption}  % 用于自定义图表标题
\usepackage{amsmath, amssymb} % 用于数学公式
\usepackage{listings} % 用于插入代码
\usepackage{xcolor}   % 用于代码高亮颜色

% 设置代码块样式
\lstset{
    basicstyle=\ttfamily\small,
    keywordstyle=\color{blue},
    commentstyle=\color{green!60!black},
    stringstyle=\color{red},
    breaklines=true,
    numbers=left,
    numberstyle=\tiny,
    frame=single,
    tabsize=4
}

% 设置图表标题格式
\captionsetup[figure]{name=图, labelsep=period}
\captionsetup[table]{name=表, labelsep=period}

% ========== 正文开始 ==========
\begin{document}

% --- 标题与作者信息 ---
\begin{center}
    {\huge \textbf{Week2 }} \\[0.5cm]
    {\large 命令行环境；开发环境与工具；Debugging and Profiling} \\
    {\large 姓名：赵均临 \quad 学号：25020007178} \\
    {\large 日期：\today}
\end{center}
\vspace{0.5cm}

% --- 目录 ---
\tableofcontents
\newpage

% --- 第一章：课上实验题目 ---
\section{课上实验题目}
\subsection{实验5：控制一个可清理的后台任务}
详细描述你所使用的理论或算法。
例如，一个行内公式：$E = mc^2$。
\begin{lstlisting}[language=Python]
import numpy as np
\%
print("Hello World")
\end{lstlisting}




\subsection{实验6：语义重构与本地开发反馈}
介绍你的实验设计或系统架构。
可以插入一张图片，例如：
\begin{figure}[H] % [H] 表示图片固定在此处
    \centering
    \includegraphics[width=0.7\textwidth]{example-image} % 请替换为你的图片文件名
    \caption{系统架构图}
    \label{fig:system}
\end{figure}

\subsection{实验7：用调试器定位归并排序缺陷}
介绍你的实验设计或系统架构。
可以插入一张图片，例如：
\begin{figure}[H] % [H] 表示图片固定在此处
    \centering
    \includegraphics[width=0.7\textwidth]{example-image} % 请替换为你的图片文件名
    \caption{系统架构图}
    \label{fig:system}
\end{figure}


\subsection{实验8：先测量，再优化慢速词频程序}
介绍你的实验设计或系统架构。
可以插入一张图片，例如：
\begin{figure}[H] % [H] 表示图片固定在此处
    \centering
    \includegraphics[width=0.7\textwidth]{example-image} % 请替换为你的图片文件名
    \caption{系统架构图}
    \label{fig:system}
\end{figure}

% --- 第二章：练习 ---
\section{练习：}
\subsection{实验环境}
介绍实验的软硬件环境。

\subsection{结果分析}
对实验结果进行分析。
也可以插入一个表格：
\begin{table}[H]
    \centering
    \begin{tabular}{|c|c|c|}
        \hline
        模型 & 准确率 & 召回率 \\
        \hline
        模型A & 95.2\% & 94.8\% \\
        模型B & 97.1\% & 96.5\% \\
        \hline
    \end{tabular}
    \caption{不同模型的性能对比}
    \label{tab:performance}
\end{table}

这是一个独立公式的例子：
\begin{equation}
    \label{eq:loss}
    \mathcal{L} = -\frac{1}{N}\sum_{i=1}^{N} \log p(y_i | x_i)
\end{equation}

% --- 第三章：练习 ---
\section{练习：}
\subsection{填充}


% --- 第四章：实验心得 ---
\section{实验心得}
总结全文，并指出未来可以改进的方向。

\end{document}